\chapter{Tech Portfolio: The Common Language of Growth\\ภาษากลางของการเติบโต}

CV บอกว่า Mentee เคยทำอะไร แต่ Tech Portfolio ช่วยให้ทั้ง Mentor และ Mentee เข้าใจว่า ประสบการณ์เหล่านั้นรวมกันเป็นคุณค่าแบบใด มีหลักฐานเพียงใด และควรเติบโตไปทางไหน จึงไม่ใช่เอกสารนิ่งเพื่อสมัครงาน แต่เป็นเครื่องมือสนทนาที่เปลี่ยนแปลงตามการเรียนรู้และบริบท

\learningobjectives{
  \item อธิบายความแตกต่างระหว่าง CV กับ Tech Portfolio
  \item ใช้องค์ประกอบทั้งเจ็ดเพื่อวินิจฉัยช่องว่างการเติบโต
  \item วางวงจรทบทวน Portfolio ที่นำไปสู่การลงมือทำ
}

\section{จากรายการผลงานสู่เรื่องราวของคุณค่า}
นักวิจัยจำนวนมากมีความเชี่ยวชาญลึก แต่คนนอกสาขากลับมองไม่เห็นคุณค่า เพราะข้อมูลกระจายอยู่ในรายชื่อบทความ ทุน และตำแหน่ง Tech Portfolio เชื่อมข้อมูลเหล่านั้นให้ตอบคำถามสามข้อ: \textbf{ท่านเป็นใครและทำอะไรได้ดี} \textbf{มีสิ่งใดยืนยัน} และ \textbf{คุณค่านั้นเปลี่ยนอะไรให้ใคร}

\begin{table}[ht]
\centering
\begin{tabularx}{\textwidth}{P{.21\textwidth}YY}
\toprule
 & \textbf{CV} & \textbf{Tech Portfolio}\\
\midrule
หน่วยหลัก & ตำแหน่งและผลงาน & คุณค่า กระบวนการ และการเรียนรู้\\
ผู้อ่าน & ผู้ประเมินคุณสมบัติ & Mentor ผู้ร่วมงาน แหล่งทุน ผู้ใช้\\
ลักษณะ & บันทึกย้อนหลัง & เครื่องมือสนทนาและออกแบบอนาคต\\
คำถามหลัก & เคยทำอะไร & ทำอย่างไร มีหลักฐานอะไร และจะไปไหนต่อ\\
\bottomrule
\end{tabularx}
\end{table}

\section{องค์ประกอบทั้งเจ็ด}
\begin{longtable}{P{.2\textwidth}P{.34\textwidth}P{.34\textwidth}}
\toprule
\textbf{องค์ประกอบ} & \textbf{คำถามหลัก} & \textbf{ความเสี่ยงเมื่อขาด}\\
\midrule
Profile & ฉันคือใครในระบบนิเวศนี้ & คนอื่นไม่เข้าใจจุดยืน\\
Skills & ฉันทำอะไรที่สร้างคุณค่าได้ & ความเชี่ยวชาญคลุมเครือ\\
Projects & ฉันแก้ปัญหาใดและคิดอย่างไร & เห็นเพียงกิจกรรม ไม่เห็นความสามารถ\\
Process & วิธีทำงานทำซ้ำและถ่ายทอดได้หรือไม่ & ผลสำเร็จพึ่งบุคคล\\
Evidence & สิ่งใดยืนยันข้อกล่าวอ้าง & ความน่าเชื่อถือต่ำ\\
Impact & ใครเปลี่ยนแปลงอย่างไร & ผลงานไม่เชื่อมคุณค่าสาธารณะ\\
Reflection & เรียนรู้อะไรและจะเปลี่ยนอย่างไร & ทำซ้ำโดยไม่พัฒนา\\
\bottomrule
\end{longtable}

องค์ประกอบทั้งเจ็ดมิได้เป็นกล่องแยกจากกัน Profile ที่น่าเชื่อถือต้องมี Skills รองรับ; Projects ทำให้ Skills มองเห็นได้; Process อธิบายความสามารถในการทำซ้ำ; Evidence ยืนยันผล; Impact แปลผลสู่คุณค่าของผู้มีส่วนได้ส่วนเสีย; และ Reflection เชื่อมอดีตกับการตัดสินใจครั้งใหม่

\section{ใช้ Portfolio เป็นเครื่องมือวินิจฉัย}
Mentor ไม่ควรรีบแก้ถ้อยคำทุกส่วน แต่ควรอ่านหา \enquote{ช่องว่างเชิงตรรกะ} ตัวอย่างเช่น Mentee กล่าวว่ามีความเชี่ยวชาญด้านการพัฒนาเทคโนโลยีเพื่อเกษตรกร แต่ Projects ทั้งหมดยังอยู่ในห้องปฏิบัติการ ช่องว่างนี้มิได้แปลว่างานไม่มีคุณค่า เพียงชี้ว่าควรออกแบบ Evidence จากผู้ใช้และขยายความเข้าใจเรื่อง Impact

\begin{examplebox}
Portfolio ฉบับแรกของนักวิจัยอาหารระบุว่า \enquote{เชี่ยวชาญ AI และมีบทความ 18 เรื่อง} หลังสนทนาด้วยกรอบทั้งเจ็ด ข้อความพัฒนาเป็น \enquote{ออกแบบระบบคัดแยกคุณภาพผลผลิตด้วย computer vision เพื่อลดการสูญเสียหลังเก็บเกี่ยว โดยต้นแบบล่าสุดลดเวลาคัดแยกของสหกรณ์นำร่องลงร้อยละ 35} ความต่างไม่ได้อยู่ที่ภาษาสวยขึ้น แต่อยู่ที่ตรรกะของคุณค่าที่ชัดขึ้น
\end{examplebox}

\section{วงจร Portfolio}
Portfolio ที่มีชีวิตควรถูกทบทวนอย่างน้อยทุกไตรมาส: รวบรวมข้อมูลใหม่ ตีความร่วมกับ Mentor เลือกช่องว่างสำคัญ ออกแบบการทดลอง และบันทึกสิ่งที่เรียนรู้ การทบทวนเช่นนี้เปลี่ยน Portfolio จากเครื่องมือประชาสัมพันธ์เป็นระบบบริหารการเติบโต

% Figure placeholder: Seven-element Tech Portfolio cycle

\begin{mentornote}
อย่าถามเพียงว่า \enquote{ส่วนใดยังเขียนไม่ครบ} ให้ถามว่า \enquote{ส่วนใดหากชัดขึ้น จะเปลี่ยนการตัดสินใจครั้งต่อไปมากที่สุด}
\end{mentornote}

\begin{keytakeaway}
Tech Portfolio เป็นภาษากลางที่เชื่อมความเชี่ยวชาญ หลักฐาน คุณค่า และทิศทางการเติบโต การเขียนเป็นเพียงผลพลอยได้; เป้าหมายแท้จริงคือการคิดและตัดสินใจได้ดีขึ้น
\end{keytakeaway}

\begin{reflection}
เลือกโครงการหนึ่งของ Mentee และเขียนข้อมูลที่มีอยู่ใต้หัวข้อทั้งเจ็ดโดยไม่แต่งเติม ทำเครื่องหมายหัวข้อที่มีเพียงข้อกล่าวอ้างแต่ยังไม่มีหลักฐาน แล้วเลือกหนึ่งช่องว่างเพื่อใช้เป็นวาระสนทนาครั้งถัดไป
\blankline\blankline
\end{reflection}

\chaptersummary{Tech Portfolio ขยายมุมมองจากประวัติผลงานไปสู่ความสามารถที่พิสูจน์ได้และคุณค่าที่สื่อสารได้ องค์ประกอบทั้งเจ็ดช่วยให้ Mentor เห็นทั้งจุดแข็ง ช่องว่าง และการทดลองที่ควรเกิดขึ้นต่อไป}
